\documentclass[letterpaper]{article}
\usepackage[utf8]{inputenc}
\usepackage[T1]{fontenc}
\usepackage[activeacute,spanish]{babel}
\usepackage[vmargin=4cm,tmargin=3cm,hmargin=2cm,letterpaper]{geometry}%
\usepackage{helvet}
\usepackage{amsmath,amsfonts,amssymb}
\usepackage{graphicx}
\usepackage{color}
\usepackage{xcolor}
\usepackage{verbatim}
\usepackage{tabls}
\usepackage{lastpage}
\usepackage{fancyhdr}
\usepackage{url}
\usepackage{listings}
%%%%%%%%%%%%%%%%%%%%%%%%%%%%%%%%%%%%%%%%%%%%%%%%%%%%%%%%%%%%%%%%%%%%%%%%%%%%%%%%%%%%%%%
\usepackage{tikz}
\usepackage{pgf}
\usepackage{pgffor}
\usepgfmodule{plot}
\usepackage{wrapfig}
\usetikzlibrary{arrows,decorations,snakes,backgrounds,fit,calc,through,scopes,positioning,automata,chains,er,fadings,calendar,matrix,mindmap,folding,patterns,petri,plothandlers,plotmarks,shadows,shapes,shapes.arrows,topaths,trees}

\lstset{% general command to set parameter(s)
%   basicstyle=\small,
  % print whole listing small
%   keywordstyle=\color{black}\bfseries\underbar,
  % underlined bold black keywords
%   identifierstyle=,
  % nothing happens
%   commentstyle=\color{white}, % white comments
%   stringstyle=\ttfamily,
  % typewriter type for strings
  showstringspaces=false}
  % no special string spaces

\pagestyle{fancy}
\color{black}
\fancyhead{}
\renewcommand{\headrule}{\hrule\vspace*{0.5mm}\rule{\linewidth}{0.8mm}}
\renewcommand{\familydefault}{\sfdefault}

\graphicspath{{./images/}}
\lhead{\includegraphics[width=2cm]{logoucr.png}}
\rhead{\includegraphics[width=3cm]{eie-text-gray-6x3cm.png}}
\chead{UNIVERSIDAD DE COSTA RICA\\FACULTAD DE INGENIERÍA\\ESCUELA DE INGENIERÍA ELÉCTRICA\\\textbf{ESTRUCTURAS ABSTRACTAS DE DATOS Y\\ ALGORITMOS PARA INGENIERÍA}\\IE-0217\\I CICLO 2014\\PROPUESTA DE I PROYECTO LIBRERÍA C++}

\lfoot{}%
\cfoot{}%
%\cfoot{\thepage\ de \pageref{LastPage}}%
\rfoot{}%

%%%%%%%%%%%%%%%%%%%%%%%%%%%%%%%%%%%%%%%%%%%%%%%%%%%%%%%%%%%%%%%%%%%%%%%%%%%%%%%%%%%%%%%%%%%%%%%%%%%%%%%%%%%%%%%
\newcommand{\uic}{blue} %user-input color
%%%%%%%%%%%%%%%%%%%%%%%%%%%%%%%%%%%%%%%%%%%%%%%%%%%%%%%%%%%%%%%%%%%%%%%%%%%%%%%%%%%%%%%%%%%%%%%%%%%%%%%%%%%%%%%%%%
\newcommand{\uim}{\_\_} %user-input marker
%%%%%%%%%%%%%%%%%%%%%%%%%%%%%%%%%%%%%%%%%%%%%%%%%%%%%%%%%%%%%%%%%%%%%%%%%%%%%%%%%%%%%%%%%%%%%%%%%%%%%%%%%%%%%%%%%%
\newcommand{\userinput}[1]{\textcolor{\uic}{\uim#1\uim}}


%%%%%%%%%%%%%%%%%%%%%%%%%%%%%%%%%%%%%%%%%%%%%%%%%%%%%%%%%%%%%%%%%%%%%%%%%%%%%%%%%%%%%%%%%%%%%%%%%%%%%%%%%%%%%%%%%%
\begin{document}\vspace*{2cm}
%%%%%%%%%%%%%%%%%%%%%%%%%%%%%%%%%%%%%%%%%%%%%%%%%%%%%%%%%%%%%%%%%%%%%%%%%%%%%%%%%%%%%%%%%%%%%%%%%%%%%%%%%%%%%%%%%%

%%%%%%%%%%%%%%%%%%%%%%%%%%%%%%%%%%%%%%%%%%%%%%%%%%%%%%%%%%%%%%%%%%%%%%%%%%%%%%%%%%%%%%%%%%%%%%%%%%%%%%%%%%%%%%%%%%
\begin{center}
\Huge
\textbf{Kinect Data Library}
\vspace*{1cm}
\end{center}

\noindent
\small\baselineskip=14pt
\textbf{Estudiante:} \textbf{Nombre del Estudiante}\\
\textbf{Carné:} \textbf{XZYWHN}\\

%%%%%%%%%%%%%%%%%%%%%%%%%%%%%%%%%%%%%%%%%%%%%%%%%%%%%%%%%%%%%%%%%%%%%%%%%%%%%%%%%%%%%%%%%%%%%%%%%%%%%%%%%%%%%%%%%%
\section{Introducción}

En las últimas décadas el aumento del material de video es Internet ha sufrido un aumento considerable, por lo cual, el algoritmo de ... es importante para ... 

%%%%%%%%%%%%%%%%%%%%%%%%%%%%%%%%%%%%%%%%%%%%%%%%%%%%%%%%%%%%%%%%%%%%%%%%%%%%%%%%%%%%%%%%%%%%%%%%%%%%%%%%%%%%%%%%%%
\section{Objetivos}

\subsection{Objetivo General}

El objetivo general consiste en la creación de una librería en C++ para la manipulación de objetos tridimensionales.\\

\subsection{Objetivos Específicos}

Los objetivos específicos son:\\

\begin{enumerate}
\item Ampliar el conocimiento del lenguaje C++ a través de la creación de una librería en este lenguaje. 
\item Desarrollar una función de comparación entre las posiciones y las trayectorias de figuras tridimensionales en la librería.
\item Ajustar una función de impresión de las coordenadas de posición de una figura tridimensional en código C++ en la librería. 
\item Realizar una función de cortes en las coordenadas de posición de una figura tridimensional en la librería.
\item Implementar una función que consiga un conjunto de coordenadas de posición que sean simétricas a otro conjunto de coordenadas dado.
\item Obtener los datos para la revisión de las funciones de la librería utilizando un kinect.
\end{enumerate}

%%%%%%%%%%%%%%%%%%%%%%%%%%%%%%%%%%%%%%%%%%%%%%%%%%%%%%%%%%%%%%%%%%%%%%%%%%%%%%%%%%%%%%%%%%%%%%%%%%%%%%%%%%%%%%%%%%
\section{Metodología}


Se realizará la investigación sobre librerías que se necesiten para la manipulación de objetos tridimensionales. Previamente, se revisarán las 
librerías encontradas, para determinar la necesidad de nuestro proyecto. 
Una vez hecho esto, se buscará implementar una librería en el lenguaje C++.\\

Lo que se realizará primero, es una investigación sobre que funciones presentan las librerías encontradas en lenguaje C++ y estudiar el código de estas,
esto con el fin de entender la utilización de los datos en estos códigos, todo esto, en caso de haber encontrado una librería en lenguaje C++ con relación a este tema.
Para lograr esto, se leerá toda la bibliografía necesaria para poder entender de la mejor manera que es lo que se necesita exactamente. Una vez
entendido algunas de las diferentes implementaciones ya existentes, se pretende crear una librería en lenguaje C++ con diferentes funciones, entre ellas,
comparación e impresión de figuras tridimensionales, así como cortes y simetrías de la misma. \\

Será necesario también, la utilización de un
kinect para obtener datos de diferentes figuras tridimensionales de tal modo que se pueda dar una revisión de la librería. Finalmente, se buscará ajustar tal librería,
de modo que sea fácil para el usuario la utilización de la misma, para que sea utilizada como una solución de las necesidades
existentes sobre objetos tridimensionales.\\



%%%%%%%%%%%%%%%%%%%%%%%%%%%%%%%%%%%%%%%%%%%%%%%%%%%%%%%%%%%%%%%%%%%%%%%%%%%%%%%%%%%%%%%%%%%%%%%%%%%%%%%%%%%%%%%%%%
\section{Referencias}

\begin{enumerate}
\item Johnson, R., \textit{Las aventuras de un linuxero}, 2da. Ed. Springer Verlag, 2016
\item ... (\_\_No referencias de wikipedia, sino, fuentes primarias\_\_)
\end{enumerate}
	
\end{document}

